\section{Dados e protocolo de avaliação}
\label{sec:dados}
O Large Movie Review Dataset do IMDb \cite{maas} tem 25.000 resenhas
rotuladas no treino oficial e outras 25.000 no teste, com 12.500 por classe
em cada divisão (\(0\) negativo, \(1\) positivo). Na limpeza, decodificamos
entidades HTML, removemos tags e normalizamos espaços.

\subsection{Divisão dos dados}
\begin{tabularx}{\linewidth}{@{}lrrX@{}}
\toprule
\tabelahead{Papel} & \tabelahead{Original} & \tabelahead{Ensaios internos} & \tabelahead{Uso}\\
\midrule
Treino       & 19.823 & 6.144 & Ajustar representação, vocabulário e pesos.\\
Validação    &  2.478 & 1.536 & Escolher hiperparâmetros e época.\\
Calibração   &  2.478 &  --   & Ajustar temperatura após a seleção original.\\
Teste/auditoria & 25.000 & 2.048 & Avaliar a configuração já escolhida.\\
\bottomrule
\end{tabularx}
\legendinha{Tabela 1. Os ensaios internos tiram treino, validação e auditoria
das 19.823 resenhas usadas no treino original. Os ensaios de sequências e de
BPE compartilham treino e validação. Cada um tem sua própria auditoria,
com 2.048 casos balanceados e reservados.}

Das 25.000 do treino oficial, 221 foram removidas na deduplicação.
As 24.779 restantes se dividem em \(19.823+2.478+2.478\). Nos ensaios
internos, usamos prefixos de até 128 palavras. Na auditoria BPE, isso
recortou 1.535 resenhas e manteve 51,24\% dos tokens por palavras. Os
prefixos couberam no teto adicional de 256 tokens.
\fonte{docs/PROTOCOLO_SEQUENCIAS.md}{Protocolo sequencial} e
\fonte{docs/PROTOCOLO_BPE.md}{protocolo BPE}.

\subsection{Comprimento das resenhas antes dos recortes}
{\small
\begin{tabularx}{\linewidth}{@{}Xrrrrrrr@{}}
\toprule
\tabelahead{Unidade} & \tabelahead{\(n\)} & \tabelahead{Mín.} & \tabelahead{P10} &
\tabelahead{Média} & \tabelahead{Mediana} & \tabelahead{P90} & \tabelahead{Máx.}\\
\midrule
\LinhasDistribuicao
\bottomrule
\end{tabularx}}
\legendinha{Tabela 2. Palavras separadas por espaços no texto limpo.
*Frases estimadas por pontuação final. \fonteatual{seminario/dados/comprimentos.json}{Distribuições completas}.}

\subsection{Análise de possíveis vazamentos de informação}
Identificamos 123 resenhas do treino que também apareciam no teste e 98
duplicatas internas com hashes SHA-256 do texto limpo com \textit{casefold}.
Reservamos os rótulos do teste à avaliação. Cada hash aparece em uma só divisão
do estudo original. O teste mantém 201 ocorrências duplicadas.
\fonte{data/manifest.json}{Manifesto}.

Depois da divisão, ajustamos o TF-IDF, o vocabulário por palavras e as fusões do
BPE usando apenas o treino. A checagem por hash cobre igualdade textual.
Paráfrases e exposição no pré-treino seguem em aberto. A análise de erros e o
Qwen reutilizam o teste em caráter exploratório.

\subsection{Dicionário de dados resumido}
\begin{tabularx}{\linewidth}{@{}lX@{}}
\toprule
\tabelahead{Campo} & \tabelahead{Significado}\\
\midrule
\texttt{id} & Identificador derivado da divisão oficial e da posição na fonte.\\
\texttt{label} & Classe binária usada na função de perda e na avaliação.\\
\texttt{clean\_text} & Texto após decodificação de HTML e normalização de espaços.\\
\texttt{word\_count} & Número de segmentos separados por espaço no texto limpo.\\
\texttt{text\_sha256} & Impressão do texto normalizado usada para verificar duplicações.\\
\bottomrule
\end{tabularx}
\legendinha{Tabela 3. A única entrada do classificador é o texto.
\fonte{DATA_DICTIONARY.md}{Dicionário completo}.}
